% === Appendix ===

\appendix

\section*{Appendix}

\vspace{0.4em}

\newcommand{\appsection}[3]{%
  \noindent\colorbox{gray!7}{\parbox{\dimexpr\linewidth-2\fboxsep}{%
    \textbf{\textcolor{cyan!70!black}{#1}\hspace{0.5em}#2}\hfill\pageref{#3}%
  }}\par\vspace{0.05em}%
}
\newcommand{\appsubsection}[2]{%
  \noindent\hspace{1em}#1\dotfill\pageref{#2}\par%
}

\begingroup
\scriptsize
\begin{center}
\textbf{\large Table of Contents}
\end{center}

\vspace{0.25em}
\hrule
\vspace{0.35em}

\appsection{A}{Reflective Code Evolution Algorithm}{app:algorithm}
\appsubsection{A.1\hspace{0.5em}Evolution Loop}{alg:evolution-loop}
\appsubsection{A.2\hspace{0.5em}Split Evaluation}{alg:split-eval}

\vspace{0.1em}
\appsection{B}{Benchmarks and Dataset Details}{app:datasets}
\appsubsection{B.1\hspace{0.5em}Dataset Splits}{app:dataset-splits}
\appsubsection{B.2\hspace{0.5em}LoCoMo}{app:dataset-locomo}
\appsubsection{B.3\hspace{0.5em}ALFWorld}{app:dataset-alfworld}
\appsubsection{B.4\hspace{0.5em}HealthBench}{app:dataset-healthbench}
\appsubsection{B.5\hspace{0.5em}PRBench}{app:dataset-prbench}

\vspace{0.1em}
\appsection{C}{Baseline Methods}{app:baselines}
\appsubsection{C.1\hspace{0.5em}No Memory}{app:baseline-no-memory}
\appsubsection{C.2\hspace{0.5em}Retrieval-Based Systems}{app:baseline-retrieval}
\appsubsection{C.3\hspace{0.5em}Self-Evolution Systems}{app:baseline-self-evolution}
\appsubsection{C.4\hspace{0.5em}Prompt-Optimizing Systems}{app:baseline-prompt}

\vspace{0.1em}
\appsection{D}{Experiment Configuration}{app:experiment-config}
\appsubsection{D.1\hspace{0.5em}Computational Cost}{app:cost}
\appsubsection{D.2\hspace{0.5em}Hyperparameters}{app:hyperparameters}
\appsubsection{D.3\hspace{0.5em}Representative Subset Selection}{app:representative-subset}

\vspace{0.1em}
\appsection{E}{Seed Programs}{app:seeds}
\appsubsection{E.1\hspace{0.5em}Vector Search}{app:seed-vector-search}
\appsubsection{E.2\hspace{0.5em}LLM Summarizer}{app:seed-llm-summarizer}
\appsubsection{E.3\hspace{0.5em}Experience Learner}{app:seed-experience-learner}

\vspace{0.1em}
\appsection{F}{Evolved Programs}{app:evolved-programs}
\appsubsection{F.1\hspace{0.5em}ALFWorld: Deterministic Action Cache}{app:evolved-alfworld}
\appsubsection{F.2\hspace{0.5em}LoCoMo: Multi-Signal Episodic Index}{app:evolved-locomo}
\appsubsection{F.3\hspace{0.5em}HealthBench and PRBench Program Sketches}{app:evolved-health-pr}
\appsubsection{F.4\hspace{0.5em}Cross-Benchmark Structural Comparison}{app:evolved-comparison}

\vspace{0.1em}
\appsection{G}{Prompt Templates and Mutation Interface}{app:prompts}
\appsubsection{G.1\hspace{0.5em}Task Agent Prompts}{app:task-agent-prompts}
\appsubsection{G.2\hspace{0.5em}Reflector Prompt}{app:reflector-prompt}
\appsubsection{G.3\hspace{0.5em}Compile-Fix Prompt}{app:compile-fix-prompt}

\vspace{0.1em}
\appsection{H}{Evaluation Protocols}{app:eval-protocols}
\appsubsection{H.1\hspace{0.5em}Token F1 (LoCoMo)}{app:metric-token-f1}
\appsubsection{H.2\hspace{0.5em}Binary Success (ALFWorld)}{app:metric-binary-success}
\appsubsection{H.3\hspace{0.5em}Rubric-Based Scoring}{app:metric-rubric}

\vspace{0.1em}
\appsection{I}{Extended Results}{app:extended-results}
\appsubsection{I.1\hspace{0.5em}Complete Per-Category Results}{app:per-category}
\appsubsection{I.2\hspace{0.5em}Multi-Seed Stability Results}{app:stability}

\vspace{0.1em}
\appsection{J}{Limitations}{app:limitations}

\vspace{0.1em}
\appsection{K}{Broader Impact}{app:broader-impact}

\vspace{0.1em}
\appsection{L}{LLM Usage}{app:llm-usage}

\vspace{0.3em}
\hrule
\endgroup

\section{Reflective Code Evolution Algorithm}
\label{app:algorithm}

\begin{algorithm}[!htbp]
\caption{\textbf{Reflective code evolution for memory-program search.}
The left panel gives the population-based evolution loop, and the right panel gives split evaluation for a candidate program.}
\label{alg:evolution}
\begin{minipage}[t]{0.47\linewidth}
\textbf{(a) Evolution Loop} \label{alg:evolution-loop}
\begin{algorithmic}[1]
\REQUIRE Seed programs $\{S_1,\ldots,S_K\}$, budget $B$, temp.\ $\tau$
\REQUIRE Episodes $D_\text{e}$, static val $D_\text{v}$, rotate pool $D_\text{r}$
\ENSURE Best program $P^*$
\STATE $\mathcal{P} \gets \emptyset$
\FOR{$i = 1$ \TO $K$}
    \STATE $(J_i, \_) \gets \textsc{Eval}(S_i, D_\text{e}, D_\text{v})$
    \STATE $\mathcal{P} \gets \mathcal{P} \cup \{(S_i, J_i)\}$
\ENDFOR
\FOR{$t = 1$ \TO $B$}
    \STATE $P_\text{par} \gets \textsc{SoftmaxSample}(\mathcal{P}, \tau)$
    \STATE $\tilde{D} \gets \textsc{Sample}(D_\text{r},\, n_r)$
    \STATE $(\_, \mathcal{R}) \gets \textsc{Eval}(P_\text{par}, D_\text{e}, \tilde{D})$
    \STATE $P_\text{mut} \gets \textsc{Mutate}(P_\text{par}, \mathcal{R})$
    \STATE $P' \gets \textsc{CompileFix}^{k}(P_\text{mut})$ \COMMENT{$\bot$ if discarded}
    \IF{$P' \neq \bot$}
        \STATE $(J', \_) \gets \textsc{Eval}(P', D_\text{e}, D_\text{v})$
        \STATE $\mathcal{P} \gets \mathcal{P} \cup \{(P', J')\}$
    \ENDIF
\ENDFOR
\RETURN $\arg\max_{(P,J) \in \mathcal{P}} J$
\end{algorithmic}
\end{minipage}%
\hfill
\begin{minipage}[t]{0.50\linewidth}
\textbf{(b) Split Evaluation} \label{alg:split-eval}
\begin{algorithmic}[1]
\REQUIRE Program $P$, task agent $\mathcal{A}$, metric $\mu$, threshold $\theta$
\REQUIRE Episodes $D_\text{e}$, eval queries $D_\text{eval}$
\ENSURE Score $J$, diagnostic cases $\mathcal{R}=(\mathcal{F},\mathcal{S})$
\STATE \textit{// Phase 1: Knowledge ingestion}
\STATE $\text{KB} \gets P.\textsc{KnowledgeBase}(\text{toolkit})$
\FOR{$d \in D_\text{e}$}
    \STATE $\text{item} \gets \mathcal{A}.\texttt{extract}(d,\; P.\textsc{InstructionKnowledgeItem})$
    \STATE $\text{KB}.\texttt{write}(\text{item},\, d)$
\ENDFOR
\STATE \textit{// Phase 2: Retrieval and scoring}
\STATE $\mathcal{F}, \mathcal{S} \gets \emptyset, \emptyset$
\FOR{$(q, y) \in D_\text{eval}$}
    \STATE $\text{query} \gets \mathcal{A}.\texttt{formulate}(q,\; P.\textsc{InstructionQuery})$
    \STATE $c \gets \text{KB}.\texttt{read}(\text{query})$
    \STATE $\hat{y} \gets \mathcal{A}.\texttt{respond}(q, c,\; P.\textsc{InstructionResponse})$
    \STATE $s \gets \mu(\hat{y},\, y)$
    \IF{$s < \theta$}
        \STATE $\mathcal{F} \gets \mathcal{F} \cup \{(q, y, \hat{y}, c, s)\}$
    \ELSE
        \STATE $\mathcal{S} \gets \mathcal{S} \cup \{(q, y, \hat{y}, c, s)\}$
    \ENDIF
\ENDFOR
\RETURN $\bigl(\tfrac{1}{|D_\text{eval}|}\sum s,\;\; (\mathcal{F}, \mathcal{S})\bigr)$
\end{algorithmic}
\end{minipage}
\end{algorithm}

\FloatBarrier

\section{Benchmarks and Dataset Details}
\label{app:datasets}

\subsection{Dataset Splits}
\label{app:dataset-splits}

\autoref{tab:datasets} summarizes the data splits for each benchmark configuration.

\begin{table}[!htbp]
\centering
\caption{\textbf{Dataset split sizes.}
Rows report train, validation, and test sizes for every benchmark split used by the evaluation pipeline.}
\label{tab:datasets}
\small
\begin{tabular}{@{}llrrr@{}}
\toprule
\textbf{Benchmark} & \textbf{Split} & \textbf{Train} & \textbf{Val} & \textbf{Test} \\
\midrule
LoCoMo         & ---               & 272 sessions  & 1{,}440 QA & 100 QA \\
\midrule
ALFWorld       & Unseen            & 4{,}652 trajectories & 145 ep.    & 50 ep. \\
               & Seen              & 4{,}652 trajectories & 150 ep.    & 50 ep. \\
\midrule
HealthBench    & Data Tasks        & 257           & 120        & 100 \\
               & Emergency         & 260           & 122        & 100 \\
\midrule
PRBench        & Legal             & 100           & 100        & 50 \\
               & Finance           & 120           & 130        & 50 \\
\bottomrule
\end{tabular}
\end{table}

For LoCoMo, we exclude category~5 (adversarial and unanswerable questions), which tests refusal capability rather than memory retrieval, retaining 1{,}540 QA pairs across four categories~\citep{locomo}.
We reserve 100 QA pairs for the held-out test split and use the remaining 1{,}440 QA pairs as validation data.
For PRBench, hard items (\texttt{finance\_hard} and \texttt{legal\_hard}) are excluded before train, validation, and held-out test construction to avoid distortion of the fitness signal~\citep{prbench}.

\subsection{LoCoMo}
\label{app:dataset-locomo}

LoCoMo~\citep{locomo} is a multi-session conversational question answering benchmark in which the agent must recover relevant facts from long dialogue histories.
The dataset contains extended dialogues with timestamped sessions, and the questions cover single-hop recall, temporal reasoning, causal reasoning, and cross-session multi-hop reasoning.
We use the provided conversation sessions as episodes and report token-level F1 together with an LLM-judge score.

\subsection{ALFWorld}
\label{app:dataset-alfworld}

ALFWorld~\citep{alfworld} evaluates embodied task completion in simulated household environments through text interaction.
The agent must complete goals such as finding, heating, cleaning, and placing objects.
We use expert trajectories from the training split as episodes and report success rate within a 50-step budget on both seen and unseen environment splits.

\subsection{HealthBench}
\label{app:dataset-healthbench}

HealthBench~\citep{healthbench} is a medical question answering benchmark with professional rubric-based evaluation.
Each sample includes a multi-turn clinical dialogue, an ideal completion, and structured pass or fail criteria written by medical professionals.
We evaluate two categories: health data interpretation and emergency referral recognition.
Because HealthBench does not provide a training split, we construct episode pools from a held-out portion of the benchmark data and keep evaluation queries disjoint from episode construction.
We filter by category first, then use 54\% of the category data for episode construction and reserve 100 examples for held-out testing.
We report rubric score, computed as the fraction of criteria marked positive by an LLM judge.

\subsection{PRBench}
\label{app:dataset-prbench}
\looseness=-1
PRBench~\citep{prbench} evaluates professional reasoning in legal analysis and financial valuation.
Each sample contains a task prompt, an expert reference, and an importance-weighted rubric.
Similar to HealthBench, PRBench does not provide a training split, so we construct episodes from a held-out portion and evaluate on disjoint query sets.
We filter out the hard subsets before splitting, use 40\% of the remaining data for episode construction, and reserve 50 examples for held-out testing.
For consistency, we apply the same rubric-based scoring protocol used for HealthBench to both PRBench splits in our evaluation.

\FloatBarrier

\section{Baseline Methods}
\label{app:baselines}

We compare \sysname{} against a no-memory control and eight memory-based baselines spanning retrieval-based systems, self-evolution systems, and prompt-optimizing systems.

\subsection{No Memory}
\label{app:baseline-no-memory}

The no-memory control answers each task directly without any external memory store or retrieved context during evaluation.

\subsection{Retrieval-Based Systems}
\label{app:baseline-retrieval}

These methods store raw observations and retrieve relevant context by similarity.
\textit{Vector Search} stores each observation in a vector collection and retrieves top-$k$ items by embedding cosine similarity, without task-specific parsing~\citep{rag}.
\textit{G-Memory}~\citep{alma} augments vector retrieval with a hierarchical graph over related tasks implemented with SQLite.
\textit{Mem0}~\citep{mem0} extracts atomic facts from observations and maintains them with LLM-driven add, update, and delete operations as new observations arrive.

\subsection{Self-Evolution Systems}
\label{app:baseline-self-evolution}

These methods distill reusable knowledge from experience instead of storing raw episodes directly.
\textit{Trajectory Retrieval}~\citep{alma} retrieves full episode trajectories by embedding similarity.
\textit{ReasoningBank}~\citep{ouyang2025reasoningbank} extracts key insights from each episode and retrieves them with similar task descriptions.
\textit{Dynamic Cheatsheet}~\citep{alma} maintains a single global summary that is iteratively rewritten after each new experience.

\subsection{Prompt-Optimizing Systems}
\label{app:baseline-prompt}

These methods improve agent behavior by evolving the system prompt.
\textit{GEPA}~\citep{gepa} optimizes prompts through reflective mutation with Pareto-based candidate selection.
We evaluate two variants: \textit{GEPA} and \textit{GEPA + Vector Search}.
The second variant adds a vector database to GEPA because the original GEPA setup does not include a persistent knowledge base, which can limit performance on memory-intensive tasks.

\section{Experiment Configuration}
\label{app:experiment-config}

\subsection{Computational Cost}
\label{app:cost}

\autoref{tab:cost} reports the computational cost of running \sysname{} evolution and baselines across all benchmarks.
We break down API calls by role: the \emph{task agent} (knowledge extraction, query generation, and answer generation), the \emph{reflector} (code mutation via LLM reflection), the \emph{judge} (per-criterion rubric scoring for HealthBench and PRBench), and the \emph{toolkit} (any LLM call issued from inside the memory pipeline at write or read time, including evolved \texttt{write()}/\texttt{read()} methods, baseline retrieval steps, and embedding-driven preprocessing helpers).
All experiments use \texttt{gpt-5.4-mini} as the task agent, judge, and toolkit model, and \texttt{gpt-5.3-codex} as the reflector model~\citep{openai2026gpt54mini,openai2026gpt53codex}.
Cost estimates use Azure OpenAI pricing: \$0.75/\$4.50 per 1M input/output tokens for \texttt{gpt-5.4-mini}, and \$1.75/\$14.00 per 1M input/output tokens for \texttt{gpt-5.3-codex}~\citep{openai2026gpt54mini,openai2026gpt53codex}.
Note that \texttt{gpt-5.3-codex} is a reasoning model; its output token counts include internal reasoning tokens, which are billed at the same rate as visible output tokens when computing the estimates below.
Token columns report GPT-family LLM tokens.
BGE-M3 embedding tokens used for vector search and subset selection are tracked separately from the GPT token and dollar estimates.

\begin{table}[!htbp]
\centering
\caption{\textbf{Computational cost of evolution runs and representative baselines.}
Rows report iteration count, API calls by role, GPT-family token counts, estimated dollar cost, and wall-clock time.
Wall-clock time includes environment interaction overhead for ALFWorld runs.}
\label{tab:cost}
\small
\setlength{\tabcolsep}{4pt}
\begin{tabular}{@{}ll r rrrr rr r r@{}}
\toprule
\textbf{Benchmark} & \textbf{Config} & \textbf{Iter.} & \textbf{Task} & \textbf{Reflect.} & \textbf{Judge} & \textbf{Toolkit} & \textbf{Input} & \textbf{Output} & \textbf{Cost} & \textbf{Time} \\
\midrule
LoCoMo        & \sysname{} & 20 & 5802 & 48 & --- & 135 & 4.9M & 4.2M & \$25.43 & 5.1h \\
              & Vector Search     &  0 &  200 &  0 & --- & 272 & 294K &  71K &  \$0.54 &  2m \\
              & No Memory       &  0 &  200 &  0 & --- & 272 & 219K &  64K &  \$0.45 &  1m \\
\midrule
ALFWorld      & \sysname{} & 20 & 3355 & 25 & --- & 443 & 7.5M & 1.4M & \$13.92 & 100h \\
(Unseen)      & Vector Search     &  0 &   80 &  0 & --- & 497 & 711K &  50K &  \$0.76 & 14m \\
              & No Memory       &  0 &   80 &  0 & --- & 497 & 708K &  50K &  \$0.76 & 11m \\
\midrule
ALFWorld      & \sysname{} & 20 & 2420 & 28 & --- & 494 & 5.6M & 1.1M & \$11.94 & 101h \\
(Seen)        & Vector Search     &  0 &   80 &  0 & --- & 497 & 692K &  49K &  \$0.74 & 11m \\
              & No Memory       &  0 &   50 &  0 & --- & 224 & 264K &  25K &  \$0.31 &  0m \\
\midrule
HB-Emergency  & \sysname{} & 20 & 4272 & 31 & 21370 & 2219 & 21M & 6.1M & \$47.20 & 4.3h \\
              & Vector Search     &  0 &   45 &  0 &  1102 &  423 & 830K & 145K &  \$1.28 &  6m \\
              & No Memory       &  0 &   54 &  0 &  1102 &  412 & 761K & 137K &  \$1.19 &  1m \\
\midrule
HB-DataTasks  & \sysname{} & 20 & 4204 & 24 & 15006 & 1575 & 18M & 6.3M & \$44.64 & 2.1h \\
              & Vector Search     &  0 &   24 &  0 &  1097 &  433 & 950K & 177K &  \$1.51 &  7m \\
              & No Memory       &  0 &   34 &  0 &  1105 &  423 & 836K & 171K &  \$1.40 &  2m \\
\midrule
PR-Legal      & \sysname{} & 20 & 3065 & 28 & 19984 & 1242 & 48M & 11M & \$90.19 & 15h \\
              & Vector Search     &  0 &   15 &  0 &   958 &  185 & 818K & 130K &  \$1.20 &  6m \\
              & No Memory       &  0 &   15 &  0 &   950 &  185 & 712K & 117K &  \$1.06 &  1m \\
\midrule
PR-Finance    & \sysname{} & 20 & 1795 & 31 & 18602 & 2497 & 47M & 10M & \$84.81 & 2.6h \\
              & Vector Search     &  0 &    8 &  0 &   845 &  212 & 740K & 153K &  \$1.24 &  6m \\
              & No Memory       &  0 &   13 &  0 &   895 &  207 & 646K & 139K &  \$1.11 & 11m \\
\bottomrule
\end{tabular}
\vspace{0.3em}

{\footnotesize Call counts include recorded API attempts.
Token totals sum provider usage fields when available; failed attempts without usage metadata contribute to call counts only.}
\end{table}

The total evolution cost ranges from \$12--90 per benchmark over 20 iterations, depending primarily on the evaluation protocol.
Rubric-graded benchmarks (HealthBench and PRBench) are substantially more expensive because each validation item requires multiple LLM judge calls for per-criterion scoring, making the judge the dominant cost component (60--80\% of total calls).
Token-F1 benchmarks (LoCoMo) and binary-success benchmarks (ALFWorld) avoid judge calls entirely, keeping evolution under \$26.
ALFWorld's long wall-clock time (100h) is dominated by TextWorld environment interaction; API latency contributes less.
Although the reflector makes only 24--48 calls per run, it uses the more expensive \texttt{gpt-5.3-codex} model (\$14/1M output tokens vs.\ \$4.50 for \texttt{gpt-5.4-mini}) and produces long code patches, so it accounts for 7--24\% of total cost depending on the benchmark.

Compared to baselines, the evolution overhead is a one-time cost: once the best program is found, inference uses the same number of calls as any baseline.
The per-query inference cost of an evolved program equals that of the corresponding baseline configuration (No Memory or Vector Search) plus any toolkit calls the evolved program makes, typically 0--2 additional LLM calls per query.


% Stability table moved to §4 Results (results.tex)

\FloatBarrier

\subsection{Hyperparameters}
\label{app:hyperparameters}

\autoref{tab:hyperparams} lists the key hyperparameters used across all experiments.

\begin{table}[!htbp]
\centering
\caption{\textbf{Hyperparameters for \sysname{} evolution.}
Rows list shared settings and benchmark-specific ranges used in the main experiments.}
\label{tab:hyperparams}
\small
\begin{tabular}{@{}ll@{}}
\toprule
\textbf{Parameter} & \textbf{Value} \\
\midrule
Evolution iterations         & 20 \\
Population seed programs     & 3 \\
Selection strategy           & Softmax ($\tau=0.15$) \\
Static validation subset     & 32--60 items ($k$-means clustering) \\
Rotating validation subset   & 5 items per iteration \\
Train-to-validation ratio    & 2 (train episodes = $2 \times$ validation items) \\
Embedding model (subset selection) & BGE-M3~\citep{bge_m3} \\
Reflector model              & GPT-5.3-Codex (reasoning: medium) \\
Task agent / toolkit model   & GPT-5.4-mini \\
Toolkit LLM call budget      & 1 call per \texttt{read}/\texttt{write} invocation \\
Compile-fix attempts ($k$)   & 3 \\
Evaluator thread pool        & 64 workers \\
LLM retry attempts           & 3 (exponential backoff) \\
\bottomrule
\end{tabular}
\end{table}

\subsection{Representative Subset Selection}
\label{app:representative-subset}

We construct representative subsets to reduce evaluation cost while preserving coverage of the full validation set.
For the static validation subset, we embed all validation samples, cluster them into $n$ groups using $k$-means~\citep{macqueen1967kmeans} with a fixed seed, and select the sample closest to each cluster centroid to form the fixed static validation set.
This procedure covers diverse regions of the validation distribution more effectively than random sampling under the same evaluation budget.
For the rotating validation subset used as reflection context, we re-cluster the remaining validation pool at each iteration with a different $k$-means seed (the iteration index added to a base seed), and select the sample closest to each new cluster centroid; varying the seed perturbs the cluster boundaries, so the chosen samples differ across iterations even though both clustering and centroid selection are deterministic given a seed.
The static subset determines fitness, while the rotating subset supplies diverse examples for mutation.

For episode subset selection, we use a facility-location objective to select $M$ training episodes that are relevant to the validation samples~\citep{nemhauser1978submodular}.
Given validation samples $V$ and candidate episodes $E$, we select a subset $S \subseteq E$ of size $M$ to maximize:
\begin{equation}
    \max_{S \subseteq E, |S|=M} \sum_{v \in V} \max_{e \in S} \operatorname{sim}(v, e),
\end{equation}
where $\operatorname{sim}(\cdot, \cdot)$ denotes embedding similarity.
We use greedy facility-location selection in the same embedding space~\citep{nemhauser1978submodular}.

\FloatBarrier


% ============================================================
%  E. Seed Programs
% ============================================================

\section{Seed Programs}
\label{app:seeds}

The program pool is initialized with three structurally diverse seeds designed to cover different points in the memory-program search space.
All three share identical instruction constants but differ in storage strategy and retrieval logic.
\autoref{tab:seed-comparison} summarizes their key design choices.

\begin{table}[!htbp]
\centering
\caption{\textbf{Seed program comparison.}
Rows summarize each seed program's storage strategy, retrieval strategy, and read-time LLM use.}
\label{tab:seed-comparison}
\small
\begin{tabular}{@{}llll@{}}
\toprule
\textbf{Seed} & \textbf{Storage} & \textbf{Retrieval} & \textbf{LLM in \texttt{read}} \\
\midrule
Vector Search & ChromaDB chunks & Top-$k$ cosine similarity & No \\
LLM Summarizer & In-memory list & Concatenate + LLM summarize & Yes (1 call) \\
Experience Learner & Separate lesson/fact lists & Return all & No \\
\bottomrule
\end{tabular}
\end{table}

\subsection{Vector Search}
\label{app:seed-vector-search}

The Vector Search seed (\autoref{lst:seed-vector}) implements vanilla RAG~\citep{rag}: raw text is split into 500-character paragraph-aligned chunks and stored in a ChromaDB collection; \texttt{read()} retrieves the top-5 most similar documents by embedding cosine similarity.
No LLM call is used at retrieval time.

\subsection{LLM Summarizer}
\label{app:seed-llm-summarizer}

The LLM Summarizer seed (\autoref{lst:seed-llm}) stores raw text in an in-memory list without any preprocessing.
At retrieval time, all stored text is concatenated (up to 30{,}000 characters) and passed to the toolkit LLM alongside the query for query-focused summarization.
This seed tests whether neural synthesis at read time can compensate for the lack of structured indexing.

\subsection{Experience Learner}
\label{app:seed-experience-learner}

The Experience Learner seed (\autoref{lst:seed-experience}) extracts a general lesson and a specific fact from each observation, storing them in separate lists.
At retrieval time, it returns all stored lessons and facts (truncated to 500 characters each), ignoring the query entirely.
This seed explores whether dual-track extraction with full recall outperforms selective retrieval.

\begin{lstlisting}[style=evolvedcode, caption={Seed 1: Vector Search.}, label=lst:seed-vector]
class KnowledgeBase:
    """Vanilla RAG: store text chunks in ChromaDB, retrieve by semantic similarity."""
    def __init__(self, toolkit):
        self.toolkit = toolkit
        self.collection = toolkit.chroma.get_or_create_collection("knowledge")
        self._doc_id = 0

    def write(self, item: KnowledgeItem, raw_text: str) -> None:
        chunks = self._chunk(raw_text, max_chars=500)
        for chunk in chunks:
            self.collection.add(documents=[chunk], ids=[f"doc_{self._doc_id}"])
            self._doc_id += 1

    def read(self, query: Query) -> str:
        if self._doc_id == 0:
            return "No information stored."
        results = self.collection.query(
            query_texts=[query.query_text], n_results=min(5, self._doc_id))
        docs = results["documents"][0] if results["documents"] else []
        return "\n\n".join(docs)[:3000] if docs else "No relevant information found."
\end{lstlisting}

\begin{lstlisting}[style=evolvedcode, caption={Seed 2: LLM Summarizer.}, label=lst:seed-llm]
class KnowledgeBase:
    """LLM-powered query-focused summarization over stored raw texts."""
    def __init__(self, toolkit):
        self.toolkit = toolkit
        self.raw_texts: list[str] = []

    def write(self, item: KnowledgeItem, raw_text: str) -> None:
        self.raw_texts.append(raw_text)

    def read(self, query: Query) -> str:
        if not self.raw_texts:
            return "No information stored."
        combined = "\n\n".join(self.raw_texts)[:30000]
        messages = [{"role": "user", "content":
            f"Given the following query, summarize ONLY the relevant information "
            f"from the provided texts. Be concise and factual.\n\n"
            f"Query: {query.query_text}\n\nTexts:\n{combined}"}]
        try:
            result = self.toolkit.llm_completion(messages)
        except Exception:
            result = combined
        return result[:3000]
\end{lstlisting}

\begin{lstlisting}[style=evolvedcode, caption={Seed 3: Experience Learner.}, label=lst:seed-experience]
@dataclass
class KnowledgeItem:
    lesson_learned: str = field(
        metadata={"description": "A general lesson or pattern learned from the text"})
    fact_to_remember: str = field(
        metadata={"description": "A specific fact worth remembering from the text"})

class KnowledgeBase:
    """Experience-driven learner that stores lessons and facts, returns all on read."""
    def __init__(self, toolkit):
        self.toolkit = toolkit
        self.lessons: list[str] = []
        self.facts: list[str] = []

    def write(self, item: KnowledgeItem, raw_text: str) -> None:
        self.lessons.append(item.lesson_learned)
        self.facts.append(item.fact_to_remember)

    def read(self, query: Query) -> str:
        if not self.lessons and not self.facts:
            return "No information stored."
        lessons_text = "\n".join(self.lessons)[:500]
        facts_text = "\n".join(self.facts)[:500]
        return f"Lessons:\n{lessons_text}\n\nFacts:\n{facts_text}"[:3000]
\end{lstlisting}


\FloatBarrier

% ============================================================
%  F. Evolved Programs
% ============================================================

\section{Evolved Programs}
\label{app:evolved-programs}

We present the best evolved programs for two representative benchmarks, selected to illustrate the structural diversity described in Section~\ref{sec:results}.
Full source code for all seven benchmark configurations is included in the released experiment artifacts under each run's \texttt{programs/} directory.

\subsection{ALFWorld: Deterministic Action Cache}
\label{app:evolved-alfworld}

The best ALFWorld program (\autoref{lst:alf-write}, \autoref{lst:alf-read}) constructs a deterministic action cache using SQLite.
It stores structured fields extracted from expert demonstrations---target object, destination, required state change (clean/cool/heat/examine/toggle), action hints, and failure modes---and retrieves them via a weighted scoring function that combines token overlap with exact-match bonuses for object, location, and state.
The program includes a \texttt{\_canonical\_state()} normalizer that maps diverse surface forms (``rinse'', ``wash'' $\to$ ``clean''; ``chill'', ``cold'' $\to$ ``cool'') to canonical labels, and a fallback parser that extracts task metadata from raw text when the LLM's structured extraction is sparse.
A single LLM call at read time synthesizes the top-6 retrieved memories into concise step-by-step guidance for the task agent to follow.
This description corresponds to the best ALFWorld Unseen program.
The best ALFWorld Seen program uses the same SQLite-backed design family, but its \texttt{read()} method returns deterministic ranked guidance without a toolkit LLM call.

The key architectural properties are:
\begin{itemize}
\item \textbf{SQLite-only storage}, with zero vector retrieval (ChromaDB unused).
\item \textbf{Canonical state normalization} unifies synonymous state descriptions (6 canonical states).
\item \textbf{Keyword-based scoring} with exact-match bonuses (+6 for object, +5 for location, +5 for state) and recency tiebreaker for stable ranking.
\item \textbf{Schema}: 8 structured fields per memory item (task\_summary, task\_type, target\_object, target\_location, required\_state, action\_hint, failure\_mode, keywords).
\end{itemize}

\begin{lstlisting}[style=evolvedcode, caption={\textcolor{red}{[Revised]} ALFWorld evolved \texttt{write()} --- canonical state extraction and fallback parsing.}, label=lst:alf-write]
def _canonical_state(self, text):
    t = (text or "").lower()
    if any(k in t for k in ["rinse", "wash", "clean"]): return "clean"
    if any(k in t for k in ["cool", "chill", "cold", "refrigerat"]): return "cool"
    if any(k in t for k in ["heat", "warm", "hot"]): return "heat"
    if any(k in t for k in ["examine", "inspect", "look at"]): return "examine"
    if any(k in t for k in ["toggle", "turn on", "light"]): return "toggle"
    return ""

def write(self, item, raw_text):
    task_summary = self._clean(item.task_summary)
    task_type = self._clean(item.task_type)
    target_object = self._clean(item.target_object)
    target_location = self._clean(item.target_location)
    action_hint = self._clean(item.action_hint)
    failure_mode = self._clean(item.failure_mode)
    required_state = self._canonical_state(item.required_state) or \
        self._canonical_state(f"{task_summary} {task_type} {raw_text}")
    # Fallback: extract from raw task metadata when LLM fields are sparse
    if not target_object:
        m = re.search(r"object_target:\s*([A-Za-z0-9_]+)", raw_text)
        if m: target_object = m.group(1)
    kw = set()
    for part in [task_summary, task_type, target_object, target_location]:
        kw |= self._tokenize(part)
    for k in item.keywords: kw |= self._tokenize(k)
    self.db.execute("INSERT INTO memories (...) VALUES (...)",
        (task_summary, task_type, target_object, target_location,
         required_state, action_hint, failure_mode, json.dumps(sorted(kw)),
         raw_text[:2000]))
\end{lstlisting}

\begin{lstlisting}[style=evolvedcode, caption={\textcolor{red}{[Revised]} ALFWorld evolved \texttt{read()} --- weighted scoring with exact-match bonuses.}, label=lst:alf-read]
def read(self, query):
    rows = self.db.execute("SELECT id, task_summary, task_type, "
        "target_object, target_location, required_state, action_hint, "
        "failure_mode, keywords_json FROM memories").fetchall()
    q_obj = self._clean(query.target_object)
    q_loc = self._clean(query.target_location)
    q_state = self._canonical_state(query.required_state)
    q_tokens = self._tokenize(query.request) | self._tokenize(q_obj) | self._tokenize(q_loc)
    scored = []
    for row in rows:
        (row_id, _, _, target_object, target_location,
         required_state, _, _, keywords_json) = row
        mem_kw = set(json.loads(keywords_json or "[]"))
        score = len(q_tokens & mem_kw)                        (*@\ding{172}@*)
        if q_obj == (target_object or ""): score += 6         (*@\ding{173}@*)
        if q_loc == (target_location or ""): score += 5       (*@\ding{174}@*)
        if q_state == (required_state or ""): score += 5      (*@\ding{175}@*)
        score += min(row_id, 1000) / 100000.0                 (*@\ding{176}@*)
        scored.append((score, row))
    selected = sorted(scored, reverse=True, key=lambda x: x[0])[:6]
    # One LLM call: synthesize top memories into step-by-step guidance
    synthesized = self.toolkit.llm_completion([...])           (*@\ding{177}@*)
    return synthesized[:3000]
\end{lstlisting}
\vspace{-0.5em}
{\scriptsize \ding{172}~Token overlap \quad \ding{173}~Exact object match \quad \ding{174}~Exact location match \quad \ding{175}~State match \quad \ding{176}~Recency tiebreaker \quad \ding{177}~LLM synthesis}


\subsection{LoCoMo: Multi-Signal Episodic Index}
\label{app:evolved-locomo}

The best LoCoMo program builds a hybrid memory combining SQLite and ChromaDB (\autoref{lst:loco-schema}, \autoref{lst:loco-read}).
Each observation is parsed into 7 structured metadata fields (participants, organizations, activities, key facts, relation facts, named entities) and stored in both a relational table and a vector collection.
At retrieval time, the program fuses semantic similarity scores from ChromaDB with lexical overlap scores from SQLite, applies person-focused boosting when the query targets a specific individual, limits per-source diversity (at most 2 chunks from any single dialogue), and separately ranks extracted facts by query relevance.
The final output combines top-ranked candidate facts and relevant text excerpts for answer generation.

The key architectural properties are:
\begin{itemize}
\item \textbf{425 source lines}, with the largest hybrid retrieval procedure among the best programs.
\item \textbf{Dual storage}: SQLite for structured metadata + ChromaDB for semantic search.
\item \textbf{7 metadata fields} per item including relation facts (subject-verb-object triples).
\item \textbf{Source diversity cap}: maximum 2 chunks per source document, preventing any single conversation from dominating retrieval for a query.
\item \textbf{Two-tier output}: candidate facts (direct answer candidates) and relevant excerpts (contextual evidence for grounded answer generation in responses).
\end{itemize}

\begin{lstlisting}[style=evolvedcode, caption={LoCoMo evolved schema --- 7 structured metadata fields.}, label=lst:loco-schema]
@dataclass
class KnowledgeItem:
    summary: str = field(metadata={"description": "Short summary"})
    participants: list[str] = field(
        metadata={"description": "People explicitly discussed in the text"})
    organizations: list[str] = field(
        metadata={"description": "Organizations, groups, or institutions"})
    activities: list[str] = field(
        metadata={"description": "Activities, hobbies, events mentioned"})
    key_facts: list[str] = field(
        metadata={"description": "Atomic factual statements useful for QA"})
    relation_facts: list[str] = field(
        metadata={"description": "Subject-verb-object facts"})
    named_entities: list[str] = field(
        metadata={"description": "Exact named items: orgs, games, places"})
\end{lstlisting}

\begin{lstlisting}[style=evolvedcode, caption={\textcolor{red}{[Revised]} LoCoMo evolved \texttt{read()} excerpt --- hybrid scoring and source diversity.}, label=lst:loco-read]
def read(self, query):
    # Phase 1: Semantic retrieval from ChromaDB
    results = self.collection.query(query_texts=[query_text], n_results=24)
    ids = results["ids"][0]; dists = results["distances"][0]
    for rank, (doc_id, dist) in enumerate(zip(ids, dists)):
        semantic_scores[doc_id] = 1.25 / (rank + 1) + 1.0 / (1 + dist)

    # Phase 2: Lexical scoring from SQLite
    for row in rows:
        lexical = self._overlap_score(query_tokens, doc_tokens)
        person_boost = 1.2 if focus_person in participants else 0.0
        scores[doc_id] += 2.1 * lexical + 1.0 * info_overlap + person_boost

    # Phase 3: Source diversity enforcement
    source_counts = collections.Counter()
    for doc_id, _ in ranked_pairs:
        source_id = row_by_id[doc_id]["source_id"]
        if source_counts[source_id] >= 2: continue     (*@\ding{172}@*)
        ranked_ids.append(doc_id)
        source_counts[source_id] += 1

    # Phase 4: Separate fact ranking
    for doc_id in ranked_ids:
        for fact in relation_facts + key_facts:
            fscore = overlap(query_tokens, fact_tokens)
            if focus_person in fact: fscore += 0.5      (*@\ding{173}@*)
            fact_candidates.append((fscore, fact))

    # Output: candidate facts + relevant excerpts
    return "Candidate facts:\n..." + "\nExcerpts:\n..."  (*@\ding{174}@*)
\end{lstlisting}
\vspace{-0.5em}
{\scriptsize \ding{172}~Max 2 chunks per source \quad \ding{173}~Person-focused boost \quad \ding{174}~Two-tier output format}


\subsection{HealthBench and PRBench Program Sketches}
\label{app:evolved-health-pr}

The HealthBench and PRBench programs evolve toward criterion-aware retrieval rather than broad episodic recall.
The HealthBench variants use SQLite tables keyed by clinical topic, urgency signals, patient attributes, and rubric-relevant safety notes; the Data Tasks variant also uses ChromaDB for semantic retrieval.
At read time, they prioritize memories that match the current clinical intent and then use one toolkit LLM call to synthesize concise guidance for the task agent.
This structure reflects the benchmark's grading protocol: answers must cover medically relevant criteria while avoiding unsafe or unsupported advice.

The PRBench variants emphasize domain-specific decomposition.
The legal program stores issue statements, governing rules, factual predicates, counterarguments, and citation-like anchors, then retrieves by matching the requested legal theory and factual pattern.
The finance program stores assumptions, valuation drivers, comparable-company cues, risk factors, and calculation notes, then combines structured scoring with vector retrieval for rubric-targeted evidence.
These programs illustrate the same pattern as ALFWorld and LoCoMo: each task induces a different schema and retrieval procedure.

\subsection{Cross-Benchmark Structural Comparison}
\label{app:evolved-comparison}

\autoref{tab:evolved-comparison} summarizes the architectural differences across all four benchmark domains, illustrating how reflective code evolution discovers structurally distinct solutions for each task.

\begin{table}[!htbp]
\centering
\caption{\textbf{Structural comparison of best evolved programs across benchmark configurations.}
Rows report program size, storage backends, read-time LLM use, schema size, and test score.}
\label{tab:evolved-comparison}
\small
\setlength{\tabcolsep}{3pt}
\begin{tabular}{@{}lcccccc@{}}
\toprule
\textbf{Benchmark} & \textbf{Lines} & \textbf{SQLite} & \textbf{ChromaDB} & \textbf{LLM in \texttt{read}} & \textbf{Schema fields} & \textbf{Test} \\
\midrule
LoCoMo              & 425  & Yes & Yes & No  & 7  & 0.459 \\
ALFWorld (Unseen)   & 276  & Yes & No  & Yes & 8  & 0.881 \\
ALFWorld (Seen)     & 441  & Yes & No  & No  & 8  & 0.700 \\
HB-Emergency        & 393  & Yes & No  & Yes & 7  & 0.493 \\
HB-Data             & 769  & Yes & Yes & Yes & 12 & 0.390 \\
PR-Legal            & 1184 & Yes & Yes & Yes & 22 & 0.660 \\
PR-Finance          & 509  & Yes & Yes & Yes & 10 & 0.586 \\
\bottomrule
\end{tabular}
\vspace{0.3em}

{\footnotesize Lines count the released best-program source files.}
\end{table}

\FloatBarrier


% ============================================================
%  G. Prompt Templates
% ============================================================

\section{Prompt Templates and Mutation Interface}
\label{app:prompts}

\looseness=-1
This section documents the key LLM prompts used in \sysname{}.
We categorize prompts by their role in the system: task agent prompts (used during evaluation) and reflector prompts (used during evolution).

\subsection{Task Agent Prompts}
\label{app:task-agent-prompts}

The task agent interacts with evolved memory programs through three fixed prompt templates.
These templates are parameterized by the program's \texttt{INSTRUCTION\_*} constants, meaning evolution can steer the agent's behavior by modifying these constants without changing the prompt structure.

\paragraph{Knowledge item generation.}
Given raw text and a \texttt{KnowledgeItem} schema, the task agent extracts structured information for a \texttt{write} call:
\begin{lstlisting}[style=evolvedcode, basicstyle=\ttfamily\scriptsize]
{INSTRUCTION_KNOWLEDGE_ITEM}

Text: {raw_text}

The KnowledgeItem must conform to this schema:
{schema}

Output ONLY a valid JSON object matching the schema fields. No explanation.
\end{lstlisting}

\paragraph{Query generation.}
Given a question and a \texttt{Query} schema, the task agent formulates a retrieval query for a \texttt{read} call:
\begin{lstlisting}[style=evolvedcode, basicstyle=\ttfamily\scriptsize]
{INSTRUCTION_QUERY}

Question: {question}

The query must be a JSON object matching this schema:
{schema}

Respond with the JSON only.
\end{lstlisting}

\paragraph{Answer generation.}
Given retrieved memory, the task agent generates an answer.
The \texttt{ALWAYS\_ON\_KNOWLEDGE} constant, if non-empty, is prepended to the retrieved content:
\begin{lstlisting}[style=evolvedcode, basicstyle=\ttfamily\scriptsize]
<retrieved_memory>
{ALWAYS_ON_KNOWLEDGE}

{retrieved}
</retrieved_memory>

{INSTRUCTION_RESPONSE}
\end{lstlisting}

\subsection{Reflector Prompt}
\label{app:reflector-prompt}

The reflector receives the current program, its evaluation score, underperforming cases, and optional context (reference programs, lineage history), and outputs a V4A patch to improve the program.
The complete prompt template is shown below.
Sections in \texttt{\{braces\}} are populated dynamically at each iteration; optional sections (lineage log, write examples, success cases, reference programs) are included only when available for that iteration.

\paragraph{Interface specification.}
The following specification is provided to the reflector to define the \texttt{KnowledgeBase} API:

\begin{lstlisting}[style=evolvedcode, basicstyle=\ttfamily\scriptsize]
You are designing a Knowledge Base Program that implements three classes:

1. **KnowledgeItem** (dataclass): Defines what information is captured
   as knowledge items when writing to the knowledge base.
   - Must be a @dataclass with typed fields
   - An external LLM will populate instances by generating JSON matching
     your field definitions
   - **Field types MUST be JSON-compatible**: use only str, int, float,
     bool, list[str], Optional[str]
   - Do NOT use datetime, tuple, bytes, or custom objects
   - Use `field(metadata={"description": "..."})` to describe fields

2. **Query** (dataclass): Defines what parameters are used when reading
   from the knowledge base.
   - Same constraints as KnowledgeItem

3. **KnowledgeBase** (class): The core knowledge base system.
   - `__init__(self, toolkit)`: Receives a Toolkit with:
     - `toolkit.db`: sqlite3.Connection (in-memory SQLite)
     - `toolkit.chroma`: chromadb ephemeral client
     - `toolkit.llm_completion(messages, **kwargs) -> str`: LLM for
       reasoning, summarization, and information extraction
       (1 call per write/read invocation)
     - `toolkit.logger.debug(message)`: Debug logging
   - `write(self, item: KnowledgeItem, raw_text: str) -> None`
   - `read(self, query: Query) -> str`

Allowed imports: json, re, math, hashlib, collections, dataclasses,
typing, datetime, textwrap, sqlite3, chromadb

## Runtime Constraints
- `read()` output limit: at most 3000 characters.
- `write()` / `read()` timeout: 60 seconds each.
- `toolkit.llm_completion()` budget: at most 1 LLM call per
  `write()` or `read()` invocation. The budget resets before each call.

## Instruction Constants (required)
Four module-level string constants:
- INSTRUCTION_KNOWLEDGE_ITEM: What to extract and how to structure it.
- INSTRUCTION_QUERY: How to formulate retrieval queries.
- INSTRUCTION_RESPONSE: Answer format, length, and style.
- ALWAYS_ON_KNOWLEDGE: Persistent context injected into every task
  agent prompt. Can be empty.
\end{lstlisting}

\paragraph{Patch format specification.}
The reflector is instructed to output changes in V4A patch format:

\begin{lstlisting}[style=evolvedcode, basicstyle=\ttfamily\scriptsize]
Before the patch, output a commit message summarizing your changes:

*** Commit Message
Title: <one-line summary of what you changed and why>
- <root cause / diagnosis>
- <what you changed>

Then output your changes as a V4A patch.

IMPORTANT: You MUST output the exact markers `*** Begin Patch` and
`*** End Patch` on their own lines. Do NOT wrap them in code fences.

Format:
*** Begin Patch
*** Update File: program.py
@@ <optional context hint>
 context line (1-2 lines before change)
-removed line
+added line
 context line (1-2 lines after change)
*** End Patch

Rules:
- Lines prefixed with `-` are removed, `+` are added,
  ` ` (space) are unchanged context.
- Include 1-2 context lines before and after each change.
- Multiple hunks are allowed within one `*** Update File` block.
\end{lstlisting}

\paragraph{Main reflection prompt.}
The following is the complete prompt template sent to the reflector LLM for each mutation step in evolution.

\begin{lstlisting}[style=evolvedcode, basicstyle=\ttfamily\scriptsize]
You are an expert Python programmer specializing in knowledge base
system design.

Your task: Given a Knowledge Base Program, its evaluation score, and
underperforming cases, identify the root cause of each low score and
improve the program. Improvements are two-fold:
(A) **Prompt Optimization** -- tune the four instruction constants
    (especially ALWAYS_ON_KNOWLEDGE) to steer the task agent's
    behavior, and
(B) **Memory Design** -- improve the KnowledgeItem/Query schemas and
    KnowledgeBase storage/retrieval logic.
Both dimensions matter and should be considered together.

<interface_spec>
{KB_INTERFACE_SPEC}
</interface_spec>

<rules>
1. Output your diagnosis first, then your changes as a patch.
2. The code must define exactly three classes (KnowledgeItem, Query,
   KnowledgeBase) and four module-level string constants
   (INSTRUCTION_KNOWLEDGE_ITEM, INSTRUCTION_QUERY,
   INSTRUCTION_RESPONSE, ALWAYS_ON_KNOWLEDGE).
3. KnowledgeBase.__init__ must accept `toolkit`; write takes a
   KnowledgeItem and a raw_text string; read takes a Query and
   returns str.
4. `read()` must return at most 3000 characters.
5. Keep it simple. Make minimal changes that generalize beyond the
   specific cases shown -- no hardcoded word lists or
   case-specific pattern rules.
6. **Prompt Optimization**: Update INSTRUCTION_* to steer the task
   LLM's output format. Update ALWAYS_ON_KNOWLEDGE with domain
   strategies, heuristics, and behavioral rules the task agent
   should always follow -- this constant is injected into EVERY
   task agent action/decision prompt and is often the
   highest-leverage change. Study the <model_generation> transcripts
   in the underperforming cases to identify agent behavioral
   patterns (e.g., looping, inefficient exploration, wrong object
   selection) that ALWAYS_ON_KNOWLEDGE can fix.
7. **Memory Design**: Improve KnowledgeItem/Query field schemas and
   KnowledgeBase read()/write() logic to store and retrieve more
   useful information for the task agent.
8. Add clear comments explaining WHY each part of the code works
   the way it does -- this helps future iterations understand and
   preserve your design decisions.
</rules>

<patch_format>
{PATCH_FORMAT_SPEC}
</patch_format>

<current_program iteration="{iteration}">
```python
{code}
```
</current_program>

<evaluation_score>{score}</evaluation_score>

{lineage_section}
{train_section}
{memory_debug_logs}
{success_section}
{reference_section}

The following cases show poor performance on the validation set after
memory has been written. Each case contains the full
retrieval-and-answer conversation trajectory.

<underperforming_cases>
<case id="1">
<question>{question}</question>
<rationale>{expected_answer}</rationale>
<model_generation>{model_output}</model_generation>
<score>{case_score}</score>
<conversation>
  [user]: {query_generation_prompt}
  [assistant]: {query_json}
  [user]: <retrieved_memory>...</retrieved_memory> {instruction}
  [assistant]: {answer}
</conversation>
</case>
...
</underperforming_cases>

<task>
1. Diagnose why these cases scored low -- examine both the retrieval
   conversation AND the <model_generation> transcript for agent
   behavioral issues.
2. Propose improvements along two dimensions:
   (A) **Prompt Optimization**: How should INSTRUCTION_* and
       ALWAYS_ON_KNOWLEDGE change to steer the task agent better?
   (B) **Memory Design**: How should the schemas or
       storage/retrieval logic change to provide more useful
       information?
3. Output your changes as a patch.
</task>
\end{lstlisting}

\paragraph{Optional sections.}
The following sections are conditionally included when their data is available:

\begin{itemize}
\item \textbf{Lineage log}: Evolution history of the current program's lineage (ancestors, children, regression markers), formatted as commit-style entries with delta scores.
Regression markers (\texttt{$\leftarrow$ REGRESSION}) flag changes that hurt performance, instructing the reflector not to repeat them.

\item \textbf{Write examples}: Sample knowledge ingestion trajectories showing how the external LLM generates knowledge items from raw document text and how \texttt{write()} is called.

\item \textbf{Success cases}: Cases where the current program performed well, instructing the reflector to preserve the behavior that makes these work.

\item \textbf{Reference programs}: Higher- or lower-scoring programs from the population, with instructions to study which design patterns (e.g., use of \texttt{llm\_completion}, ChromaDB vs.\ SQLite, schema granularity) correlate with scores within the pool.

\item \textbf{Memory debug logs}: Outputs of \texttt{toolkit.logger.debug()} calls within \texttt{write()} and \texttt{read()}, providing visibility into program execution.
\end{itemize}

\paragraph{Underperforming case selection.}
At each iteration, 2 failed cases are sampled from the rotating validation set using the Efraimidis--Spirakis weighted sampling algorithm with weight $w_i = 1 - \text{score}_i$, biasing selection toward lower-scoring cases while maintaining diversity~\citep{efraimidis2006weighted}.

\subsection{Compile-Fix Prompt}
\label{app:compile-fix-prompt}

When a mutated program fails to compile or violates runtime constraints, a compile-fix prompt is sent to the reflector with the failing program and error details:
\begin{lstlisting}[style=evolvedcode, basicstyle=\ttfamily\scriptsize]
You are an expert Python programmer. A Knowledge Base Program failed
to compile or run. Fix the error and output your fix as a patch.

{KB_INTERFACE_SPEC}

## Failing Code
```python
{code}
```

## Error
**{error_type}**: {error_details}

Fix the error and output your fix as a patch.
\end{lstlisting}
Up to $k$ compile-fix iterations are attempted before the mutation is discarded.


\FloatBarrier

% ============================================================
%  H. Evaluation Protocols
% ============================================================

\section{Evaluation Protocols}
\label{app:eval-protocols}

Each benchmark uses a task-specific evaluation metric.
All metrics return a score in $[0, 1]$; the fitness function $J(P)$ averages these scores across the evaluation set.

\subsection{Token F1 (LoCoMo)}
\label{app:metric-token-f1}

Following \citet{squad}, we compute token-level F1 between the model output and the expected answer.
Both strings are lowercased, articles (``a'', ``an'', ``the'') are removed, and all punctuation is stripped.
Precision and recall are computed over the resulting token multisets using the standard formula below:
\begin{equation}
F_1 = \frac{2 \cdot |\text{out} \cap \text{exp}|}{|\text{out}| + |\text{exp}|}.
\end{equation}
If both token sets are empty, the score is 1.0; if exactly one is empty, the score is 0.0.

\paragraph{LLM-Judge Score (LoCoMo).}
The L-J column in Table~\ref{tab:main-results} reports a complementary LLM-as-judge metric that allows for paraphrastic answers that token F1 would penalize.
A separate judge call grades each prediction against the gold answer and returns a binary correctness label, which we then average over the test set.
The judge uses the same task agent model (GPT-5.4-mini) and the prompt template below; the score is $1.0$ if the response parses as the integer $1$ and $0.0$ otherwise.
\begin{lstlisting}[style=evolvedcode, basicstyle=\ttfamily\scriptsize]
You are a strict judge. Determine if the output answers the question
correctly based on the expected answer. Reply ONLY with 1 (correct)
or 0 (incorrect).

Expected answer: <<expected>>
Actual output:   <<output>>

Score (0 or 1):
\end{lstlisting}

\subsection{Binary Success (ALFWorld)}
\label{app:metric-binary-success}

The agent interacts with the TextWorld environment for up to 50 steps~\citep{textworld,alfworld}.
A score of 1.0 is assigned if the task is completed within the step budget; 0.0 otherwise.

\subsection{Rubric-Based Scoring (HealthBench, PRBench)}
\label{app:metric-rubric}

Each evaluation item includes a structured rubric with multiple criteria, each assigned a point value (positive for desirable traits, negative for undesirable ones).
An LLM judge independently grades each criterion by examining the conversation and the last assistant response, returning a JSON object with an \texttt{explanation} and a boolean \texttt{criteria\_met} field.
For negative criteria (e.g., ``Is overly verbose''), a good response should be graded as \texttt{false} (the undesirable trait is not present).
The score is computed as:
\begin{equation}
\text{score} = \text{clip}\!\left(\frac{\sum_{i:\, \text{met}_i} p_i}{\sum_{i:\, p_i > 0} p_i},\; 0,\; 1\right),
\end{equation}
where $p_i$ is the point value for criterion $i$ and $\text{met}_i$ indicates whether the criterion was judged as met.
Negative criteria that are met subtract from the numerator.
This protocol follows the official HealthBench and PRBench evaluation methodology~\citep{healthbench,prbench}.
The LLM judge uses the following prompt template for each criterion during grading:

\begin{lstlisting}[style=evolvedcode, basicstyle=\ttfamily\scriptsize]
Your job is to look at a conversation and a rubric item, and score
the last turn (the last assistant response) on how well it follows
the rubric item.

# Conversation
<<conversation>>

# Rubric item
<<rubric_item>>

# Instructions
Return a json object with fields: "explanation" and "criteria_met".
- "explanation": why the response does or does not meet the criteria.
- "criteria_met": boolean. If a rubric item has multiple criteria,
  all must be met to return true.
- Exception: if criteria says "such as" or "for example", the
  response need not include all listed examples.
- For negative criteria (undesirable traits), return whether the
  criteria IS met (true = bad behavior present), not whether the
  response is good.
\end{lstlisting}


% ============================================================
%  I. Per-Category Results
% ============================================================

\section{Extended Results}
\label{app:extended-results}

\subsection{Complete Per-Category Results}
\label{app:per-category}

Table~\ref{tab:per-category} in the main paper reports per-category results for ALFWorld Unseen and LoCoMo with selected baselines.
Here we provide the complete per-category breakdown for all benchmarks and all baselines in the appendix tables below.

\begin{table}[!htbp]
\centering
\caption{\textbf{Complete per-category results for ALFWorld Unseen.}
Columns report success rate by task type.
Best per column appears in \textbf{bold} for each reported task category.}
\label{tab:per-cat-alf-unseen}
\small
\setlength{\tabcolsep}{4pt}
\begin{tabular}{@{}l cccccc c@{}}
\toprule
& \textbf{Simple} & \textbf{Movable} & \textbf{Clean} & \textbf{Light} & \textbf{Cool} & \textbf{Heat} & \textbf{2-Obj} \\
& \footnotesize$(n{=}8)$ & \footnotesize$(n{=}1)$ & \footnotesize$(n{=}11)$ & \footnotesize$(n{=}1)$ & \footnotesize$(n{=}7)$ & \footnotesize$(n{=}6)$ & \footnotesize$(n{=}8)$ \\
\midrule
No Memory            & 0.88 & \textbf{1.00} & 0.82 & 0.00 & 0.57 & 0.83 & 0.62 \\
Vector Search        & 0.62 & \textbf{1.00} & 0.64 & 0.00 & 0.71 & 0.67 & 0.62 \\
G-Memory             & \textbf{1.00} & \textbf{1.00} & 0.55 & \textbf{1.00} & 0.29 & 0.83 & 0.75 \\
Mem0                 & 0.88 & \textbf{1.00} & 0.64 & 0.00 & 0.57 & \textbf{1.00} & 0.75 \\
Trajectory Retrieval & 0.75 & \textbf{1.00} & 0.73 & 0.00 & 0.71 & 0.67 & 0.75 \\
ReasoningBank        & \textbf{1.00} & \textbf{1.00} & 0.55 & \textbf{1.00} & 0.57 & 0.83 & 0.75 \\
Dynamic Cheatsheet   & \textbf{1.00} & \textbf{1.00} & 0.36 & \textbf{1.00} & 0.00 & \textbf{1.00} & 0.75 \\
GEPA                 & 0.88 & \textbf{1.00} & \textbf{1.00} & 0.00 & \textbf{0.86} & \textbf{1.00} & 0.62 \\
GEPA + Vector Search & 0.88 & \textbf{1.00} & \textbf{1.00} & 0.00 & 0.57 & \textbf{1.00} & 0.62 \\
\rowcolor{green!10}
\sysname{}           & 0.88 & \textbf{1.00} & 0.91 & \textbf{1.00} & \textbf{1.00} & 0.83 & 0.75 \\
\bottomrule
\end{tabular}
\end{table}

\begin{table}[!htbp]
\centering
\caption{\textbf{Complete per-category results for ALFWorld Seen.}
Columns report success rate by task type.
Best per column appears in \textbf{bold} for each reported task category.}
\label{tab:per-cat-alf-seen}
\small
\setlength{\tabcolsep}{4pt}
\begin{tabular}{@{}l cccccc c@{}}
\toprule
& \textbf{Simple} & \textbf{Clean} & \textbf{Light} & \textbf{Cool} & \textbf{Heat} & \textbf{2-Obj} & \textbf{Overall} \\
& \footnotesize$(n{=}2)$ & \footnotesize$(n{=}12)$ & \footnotesize$(n{=}2)$ & \footnotesize$(n{=}8)$ & \footnotesize$(n{=}7)$ & \footnotesize$(n{=}15)$ & \footnotesize$(n{=}50)$ \\
\midrule
No Memory            & \textbf{1.00} & 0.33 & \textbf{1.00} & 0.50 & 0.43 & 0.87 & 0.640 \\
Vector Search        & \textbf{1.00} & 0.42 & \textbf{1.00} & \textbf{1.00} & 0.29 & 0.87 & 0.720 \\
G-Memory             & \textbf{1.00} & 0.33 & \textbf{1.00} & 0.12 & 0.29 & 0.87 & 0.560 \\
Mem0                 & \textbf{1.00} & 0.25 & \textbf{1.00} & 0.62 & 0.43 & 0.87 & 0.640 \\
Trajectory Retrieval & \textbf{1.00} & 0.58 & \textbf{1.00} & 0.88 & 0.43 & \textbf{0.93} & 0.780 \\
ReasoningBank        & \textbf{1.00} & 0.33 & \textbf{1.00} & 0.25 & 0.43 & 0.80 & 0.580 \\
Dynamic Cheatsheet   & \textbf{1.00} & 0.17 & \textbf{1.00} & 0.38 & 0.00 & 0.73 & 0.480 \\
GEPA                 & \textbf{1.00} & 0.50 & \textbf{1.00} & 0.88 & 0.43 & 0.80 & 0.720 \\
GEPA + Vector Search & \textbf{1.00} & \textbf{0.92} & \textbf{1.00} & \textbf{1.00} & 0.29 & 0.80 & \textbf{0.820} \\
\rowcolor{green!10}
\sysname{}           & \textbf{1.00} & 0.42 & \textbf{1.00} & 0.75 & \textbf{0.57} & 0.80 & 0.700 \\
\bottomrule
\end{tabular}
\vspace{0.3em}

{\footnotesize \textbf{Bold} marks the best per category; GEPA+VS has the highest overall score (0.820).}
\end{table}

\begin{table}[!htbp]
\centering
\caption{\textbf{Complete per-category results for LoCoMo.}
Columns report token F1 by question category: single-hop recall, temporal reasoning, causal reasoning, and open-domain.
Best per column appears in \textbf{bold} for each reported question category.}
\label{tab:per-cat-locomo}
\small
\setlength{\tabcolsep}{5pt}
\begin{tabular}{@{}l cccc c@{}}
\toprule
& \textbf{Single-hop} & \textbf{Temporal} & \textbf{Causal} & \textbf{Open-domain} & \textbf{Overall} \\
& \footnotesize$(n{=}18)$ & \footnotesize$(n{=}18)$ & \footnotesize$(n{=}9)$ & \footnotesize$(n{=}55)$ & \\
\midrule
No Memory            & 0.02 & 0.00 & 0.11 & 0.04 & 0.036 \\
Vector Search        & 0.19 & 0.29 & 0.38 & 0.25 & 0.256 \\
G-Memory             & 0.23 & 0.12 & 0.31 & 0.24 & 0.224 \\
Mem0                 & \textbf{0.33} & 0.26 & 0.37 & 0.43 & 0.373 \\
Trajectory Retrieval & 0.18 & \textbf{0.33} & 0.34 & 0.28 & 0.276 \\
ReasoningBank        & 0.25 & 0.07 & 0.34 & 0.19 & 0.194 \\
Dynamic Cheatsheet   & 0.14 & 0.00 & 0.22 & 0.14 & 0.124 \\
GEPA                 & 0.09 & 0.06 & 0.24 & 0.15 & 0.132 \\
GEPA + Vector Search & 0.16 & \textbf{0.38} & \textbf{0.39} & 0.31 & 0.300 \\
\rowcolor{green!10}
\sysname{}           & \textbf{0.33} & 0.19 & 0.35 & \textbf{0.51} & \textbf{0.459} \\
\bottomrule
\end{tabular}
\end{table}


HealthBench and PRBench each evaluate on single-category splits (emergency referrals, health data interpretation, legal reasoning, and financial reasoning).
Their per-category scores equal the overall scores reported in Table~\ref{tab:main-results}; we omit separate per-category tables for these benchmarks.


% ============================================================
%  J. Stability Analysis
% ============================================================

\subsection{Multi-Seed Stability Results}
\label{app:stability}

Table~\ref{tab:stability} in the main paper reports summary statistics (mean, standard deviation, coefficient of variation) across five evolution seeds.
Here we provide the full per-seed results.

\begin{table}[!htbp]
\centering
\caption{\textbf{Per-seed test scores across stability experiments.}
Best Baseline denotes the strongest fixed baseline.}
\label{tab:stability-full}
\small
\setlength{\tabcolsep}{4pt}
\begin{tabular}{@{}l ccccc ccc c@{}}
\toprule
& \multicolumn{5}{c}{\textbf{Test Score by Seed}} & & & & \\
\cmidrule(r){2-6}
\textbf{Benchmark} & \textbf{0} & \textbf{1} & \textbf{2} & \textbf{3} & \textbf{4} & \textbf{Mean} & \textbf{$\pm$Std} & \textbf{CV} & \textbf{Best Baseline} \\
\midrule
LoCoMo     & 0.459 & 0.463 & 0.364 & 0.394 & 0.423 & 0.421 & 0.042 & 10.0\% & 0.373 \\
HB-Data    & 0.390 & 0.382 & 0.357 & 0.411 & 0.413 & 0.391 & 0.023 & 5.9\% & 0.327 \\
PR-Finance & 0.586 & 0.544 & 0.582 & 0.582 & 0.513 & 0.561 & 0.032 & 5.7\% & 0.449 \\
\bottomrule
\end{tabular}
\end{table}

On LoCoMo, 4 out of 5 seeds exceed the strongest baseline (Mem0, 0.373).
The single exception (seed 2, 0.364) falls within 2.4\% of the baseline, consistent with the higher variance inherent to LoCoMo's multi-category evaluation.
On HealthBench Data and PRBench Finance, all 5 seeds exceed the respective strongest baselines (GEPA+VS at 0.327; GEPA at 0.449) by substantial margins, demonstrating consistent improvement regardless of initialization.

\autoref{tab:stability-val} reports the best validation score achieved during evolution for each seed, which determines the program selected for final test evaluation.

\begin{table}[!htbp]
\centering
\caption{\textbf{Best validation score during evolution by seed.}
Scores are computed on the static validation subset used to select the program for test evaluation.}
\label{tab:stability-val}
\small
\setlength{\tabcolsep}{4pt}
\begin{tabular}{@{}l ccccc cc@{}}
\toprule
& \multicolumn{5}{c}{\textbf{Best Validation Score}} & & \\
\cmidrule(r){2-6}
\textbf{Benchmark} & \textbf{0} & \textbf{1} & \textbf{2} & \textbf{3} & \textbf{4} & \textbf{Mean} & \textbf{$\pm$Std} \\
\midrule
LoCoMo     & 0.333 & 0.273 & 0.289 & 0.312 & 0.333 & 0.308 & 0.027 \\
HB-Data    & 0.424 & 0.381 & 0.378 & 0.397 & 0.380 & 0.392 & 0.020 \\
PR-Finance & 0.521 & 0.522 & 0.559 & 0.554 & 0.515 & 0.534 & 0.020 \\
\bottomrule
\end{tabular}
\end{table}


% ============================================================
%  J. Limitations
% ============================================================

\section{Limitations}
\label{app:limitations}

\sysname{} has three main limitations.
First, evaluating each candidate memory program requires re-ingesting episodes and re-scoring on the validation set, so the per-iteration cost grows with episode length and rubric complexity (Appendix~\ref{app:cost}); rubric-graded benchmarks reach \$90 and ALFWorld reaches 100 wall-clock hours per 20-iteration run.
Second, our experiments use a fixed task agent and reflector model pair (GPT-5.4-mini and GPT-5.3-Codex); we have not measured how strongly the discovered programs depend on this choice.
Third, evolution can drift toward task-specific heuristics encoded directly in the program logic (e.g., the synonym table in the ALFWorld action cache, Appendix~\ref{app:evolved-alfworld}), which improves in-task performance but limits transfer to other tasks (Figure~\ref{fig:task-optimized-memory}); designing search procedures that prefer reusable structure remains open.

\section{Broader Impact}
\label{app:broader-impact}

This work studies agent memory as a research object in controlled benchmark settings.
Its main positive impact is to make memory programs more reproducible and inspectable: the discovered program is executable code that can be read, tested, and released with the artifacts.
This can help researchers compare memory programs beyond fixed retrieval pipelines.
The main risk is that better memory can also make agents more persistent in carrying forward wrong, private, or biased information when the data or objective contains it.
For uses outside controlled benchmarks, we believe deployment should include data minimization, task-specific evaluation, logging, and human review, especially in high-stakes domains such as health and law.

\section{LLM Usage}
\label{app:llm-usage}

The authors used large language model agents as research assistants during the preparation of this work.
Their assistance covered early literature exploration, identification of related papers, implementation support, experiment execution and debugging, and language polishing during writing.
The authors made the scientific decisions, developed the ideas and claims, designed the method and experiments, and determined the paper structure and final content.
